\begin{center}
  \chapter{Implementation and Realization}
\end{center}
\label{chap:implementation}

\noindent
This chapter describes the concrete realization of the LiDAR Data Manager: the development environment, the implementation of each major module, and the validation of the delivered system through functional testing and performance benchmarks.

%─────────────────────────────────────────────────────────────────────────────
\section{Development Environment}

The system was developed and tested on a workstation running Ubuntu 22.04 LTS, equipped with an Intel Core i7-10700K (8 cores, 3.8\,GHz), 32\,GB DDR4 RAM, and a 1\,TB NVMe SSD. The full software stack is listed in Table~\ref{tab:software}.

\begin{table}[H]
  \centering
  \caption{Software stack and versions.}
  \label{tab:software}
  \begin{tabular}{lll}
    \toprule
    \textbf{Category} & \textbf{Tool / Library} & \textbf{Version} \\
    \midrule
    Language              & Python               & 3.12 \\
    Web Framework         & FastAPI + Uvicorn    & 0.110 / 0.29 \\
    Async MongoDB driver  & Motor                & 3.4 \\
    Database              & MongoDB Community    & 7.0 \\
    Object storage        & MinIO                & RELEASE.2024-04 \\
    Point cloud processing& PDAL                 & 2.6 \\
    Coordinate reprojection & pyproj             & 3.6 \\
    Data validation       & Pydantic v2          & 2.7 \\
    Frontend runtime      & Node.js + Vite       & 20 / 5.2 \\
    Frontend framework    & React                & 18.2 \\
    Map library           & Leaflet + react-leaflet & 1.9 / 4.2 \\
    Containerization      & Docker + Docker Compose & 26 / 2.27 \\
    \bottomrule
  \end{tabular}
\end{table}

All services (FastAPI backend, MongoDB, MinIO, and the React frontend) were orchestrated with Docker Compose, ensuring a reproducible single-command startup for both development and testing.

%─────────────────────────────────────────────────────────────────────────────
\section{LiDAR Test Data}

Two categories of datasets were used throughout development. \textbf{Synthetic data} was generated with PDAL's \texttt{readers.faux} driver, producing deterministic point clouds of known size — used primarily in the automated test suite. \textbf{Real-world data} was sourced from the IGN (Institut National de l'Information Géographique et Forestière) open data portal: airborne LiDAR tiles in LAZ format, projected in Lambert-93 (EPSG:2154), ranging from 80\,MB to 650\,MB.

\begin{table}[H]
  \centering
  \caption{Test dataset characteristics.}
  \label{tab:datasets}
  \begin{tabular}{lll}
    \toprule
    \textbf{Property} & \textbf{Synthetic} & \textbf{Real-world (IGN)} \\
    \midrule
    Format            & LAZ (LAS 1.4)        & LAZ (LAS 1.2) \\
    Point count       & 10,000 – 100,000     & 8,000,000 – 70,000,000 \\
    File size         & $<$ 1\,MB            & 80\,MB – 650\,MB \\
    Coordinate system & WGS84 (EPSG:4326)    & Lambert-93 (EPSG:2154) \\
    \bottomrule
  \end{tabular}
\end{table}

%─────────────────────────────────────────────────────────────────────────────
\section{System Interface}

\subsection{Application Overview}

The frontend is a single-page React application divided into two main views: a \textbf{dataset dashboard} for managing uploads and monitoring processing status, and a \textbf{map view} for spatial querying and zone extraction.

\begin{figure}[H]
  \centering
  \includegraphics[width=\textwidth]{figures/fig_dashboard_overview.png}
  \caption{Dataset dashboard: list of uploaded datasets with processing status badges.}
  \label{fig:dashboard}
\end{figure}

\subsection{File Upload}

Users upload LAZ or LAS files through a drag-and-drop interface. The file is streamed directly to MinIO object storage without intermediate disk writes on the server, keeping memory usage below 30\,MB regardless of file size. Once the upload completes, the dataset card appears on the dashboard with a \textit{Pending} status badge.

\begin{figure}[H]
  \centering
  \includegraphics[width=\textwidth]{figures/fig_upload_interface.png}
  \caption{Upload interface: drag-and-drop panel and upload progress indicator.}
  \label{fig:upload}
\end{figure}

\subsection{Processing Pipeline}

Clicking \textit{Process} on a dataset card triggers the backend pipeline. The status badge transitions from \textit{Pending} → \textit{Processing} → \textit{Ready}, with the frontend polling the API every two seconds.

\begin{figure}[H]
  \centering
  \includegraphics[width=0.85\textwidth]{figures/fig_processing_status.png}
  \caption{Dataset card status transitions during the processing pipeline.}
  \label{fig:status}
\end{figure}

The pipeline executes four sequential steps: metadata and bounding box extraction via PDAL, WGS84 reprojection using pyproj, flat 2D grid splitting into tiles, and COPC conversion of each tile. All intermediate and final files are stored in MinIO; tile metadata (bounding boxes, point counts, storage paths) is persisted in MongoDB.

\begin{figure}[H]
  \centering
  \includegraphics[width=\textwidth]{figures/fig_pipeline_steps.png}
  \caption{Processing pipeline: the four sequential steps executed on each uploaded file.}
  \label{fig:pipeline}
\end{figure}

\subsection{Map View and Geographic Footprint}

Once a dataset reaches \textit{Ready} status, its geographic footprint is displayed on the Leaflet map as a bounding box overlay. Users can click any dataset in the list to center the map on its extent.

\begin{figure}[H]
  \centering
  \includegraphics[width=\textwidth]{figures/fig_map_footprint.png}
  \caption{Map view: geographic footprint of a processed IGN dataset overlaid on the base map.}
  \label{fig:mapfootprint}
\end{figure}

\subsection{Zone Extraction}

Users draw a rectangular selection on the map to define a zone of interest. The frontend sends the bounding box coordinates to the \texttt{/crop} endpoint; the backend identifies intersecting tiles via a MongoDB spatial query, merges them if necessary, crops the merged point cloud to the exact bounding box using PDAL, and streams the resulting LAZ file back. The browser automatically triggers a file download.

\begin{figure}[H]
  \centering
  \includegraphics[width=\textwidth]{figures/fig_zone_selection.png}
  \caption{Zone selection: user-drawn bounding box on the map prior to extraction.}
  \label{fig:zoneselect}
\end{figure}

\begin{figure}[H]
  \centering
  \includegraphics[width=0.75\textwidth]{figures/fig_download_dialog.png}
  \caption{Browser download dialog triggered after a successful zone extraction.}
  \label{fig:download}
\end{figure}

%─────────────────────────────────────────────────────────────────────────────
\section{Key Implementation Decisions}

\subsection{Streaming Upload}

Setting \texttt{length=-1} in the MinIO SDK activates automatic multipart upload. This means the server never buffers the entire file in memory — a 650\,MB LAZ file consumes less than 30\,MB of application memory during upload.

\subsection{Dual-Schema PDAL Output Parsing}

PDAL 2.5 and 2.6 produce bounding box metadata with different JSON schemas (flat keys vs.\ nested per-dimension dictionaries). The \texttt{\_extract\_bbox} helper handles both schemas transparently, ensuring the pipeline works regardless of the PDAL version installed in the container.

\subsection{Four-Corner WGS84 Reprojection}

Reprojecting only the two extreme corners of a bounding box introduces error for projected coordinate systems like Lambert-93, where meridian convergence causes the box edges to curve. All four corners are therefore transformed, and the WGS84 envelope is taken as the min/max of the four resulting coordinates. This guarantees that the footprint displayed on the Leaflet map fully encloses the dataset.

\subsection{MongoDB Interval-Overlap Tile Query}

Spatial queries against the tile collection use the standard interval-overlap condition on both axes. A compound index on \texttt{(dataset\_id, bbox.min\_x, bbox.max\_x, bbox.min\_y, bbox.max\_y)} created at startup ensures the query is resolved via an index scan rather than a full collection scan, keeping latency below 10\,ms for collections with thousands of tile records.

%─────────────────────────────────────────────────────────────────────────────
\section{Results and Validation}

\subsection{Functional Validation}

All MVP features were delivered and validated. Table~\ref{tab:features} summarises the delivery status of each feature.

\begin{table}[H]
  \centering
  \caption{MVP feature delivery status.}
  \label{tab:features}
  \begin{tabular}{lll}
    \toprule
    \textbf{Feature} & \textbf{Status} & \textbf{Validation method} \\
    \midrule
    LAZ/LAS streaming upload        & \checkmark Complete & 650\,MB file, $<$30\,MB RAM \\
    Metadata extraction (PDAL)      & \checkmark Complete & Bounding box and SRS verified \\
    WGS84 reprojection              & \checkmark Complete & Lambert-93 footprint on map \\
    Flat 2D tiling                  & \checkmark Complete & Zero-duplication automated test \\
    COPC conversion                 & \checkmark Complete & PDAL info output verified \\
    Tile storage (MinIO + MongoDB)  & \checkmark Complete & MinIO console + mongo-express \\
    Spatial tile query              & \checkmark Complete & Correct tile subset returned \\
    Rectangular zone extraction     & \checkmark Complete & Downloaded LAZ point count verified \\
    Frontend status polling         & \checkmark Complete & Live badge updates observed \\
    3D Visualization (Potree)       & In progress         & Planned for next sprint \\
    JWT Authentication              & Planned             & Post-MVP \\
    \bottomrule
  \end{tabular}
\end{table}

\subsection{Zero-Duplication Verification}

A key correctness property of the tiling algorithm is that every point from the source file appears in exactly one tile. The automated test suite verifies this explicitly on a 10,000-point synthetic dataset for tile sizes of 100\,m, 250\,m, 500\,m, and 1000\,m. All runs produce a sum of tile point counts equal to the input point count.

\begin{figure}[H]
  \centering
  \includegraphics[width=0.85\textwidth]{figures/fig_test_output.png}
  \caption{Automated test output confirming zero duplication across 12 tiles (10,000 input points).}
  \label{fig:testoutput}
\end{figure}

\subsection{Performance Benchmarks}

Processing time was measured on three real-world IGN datasets (median of three runs each). Results are shown in Table~\ref{tab:benchmarks}.

\begin{table}[H]
  \centering
  \caption{Processing time benchmarks on development hardware.}
  \label{tab:benchmarks}
  \begin{tabular}{lrrr}
    \toprule
    \textbf{Operation} & \textbf{80\,MB (8M pts)} & \textbf{320\,MB (35M pts)} & \textbf{650\,MB (70M pts)} \\
    \midrule
    Upload to MinIO         &  4.2\,s &  17.1\,s &  34.8\,s \\
    Metadata extraction     &  8.3\,s &  31.2\,s &  63.7\,s \\
    Boundary extraction     & 12.1\,s &  47.5\,s &  96.2\,s \\
    Grid splitting          & 41.5\,s & 168.3\,s & 342.7\,s \\
    COPC conversion         & 23.8\,s &  94.1\,s & 191.3\,s \\
    Tile upload to MinIO    &  6.4\,s &  24.7\,s &  50.2\,s \\
    \midrule
    \textbf{Total}          & \textbf{96.3\,s} & \textbf{382.9\,s} & \textbf{778.9\,s} \\
    Tiles generated         & 12 & 48 & 97 \\
    Zone extraction (4 tiles) & 3.1\,s & 3.4\,s & 3.6\,s \\
    \bottomrule
  \end{tabular}
\end{table}

Processing time scales approximately linearly with point count, which is expected given that the dominant operations (grid splitting and COPC conversion) iterate over every point. Zone extraction time is effectively independent of total dataset size: it depends only on the number of tiles intersecting the query bounding box, which in these tests was held constant at four tiles.

%─────────────────────────────────────────────────────────────────────────────
\section*{Chapter Conclusion}

This chapter presented the concrete realization of the LiDAR Data Manager. The interface was demonstrated end-to-end: from file upload and pipeline processing through to geographic visualization and zone extraction. The four key implementation decisions — streaming upload, dual-schema PDAL output parsing, four-corner WGS84 reprojection, and MongoDB interval-overlap tile queries — were each motivated by a specific technical constraint encountered during development. Functional validation confirms that all MVP features have been delivered. Processing performance scales linearly with point count, and zone extraction time remains effectively constant regardless of total dataset size.